\chapter{Background and Related Works}
\label{chap:background}

\section{Retrieval-Augmented Generation}

RAG combines a generator's parametric knowledge with evidence retrieved from an
external collection \cite{lewis2020rag}. A typical pipeline transforms a user
question into a search query, retrieves a set of passages, assembles a prompt,
and asks a language model to generate an answer. The approach supports
knowledge updates without model retraining and makes source attribution
possible. It also introduces dependencies that a generator-only system does not
have: document preprocessing, index quality, search recall, ranking, context
limits, prompt construction, and grounding behavior.

The distinction between retrieval and generation is important for diagnosis.
If the required information is absent from the retrieved context, even a strong
generator cannot produce a fully grounded answer. Conversely, a relevant
context does not guarantee a faithful answer because the generator may ignore,
misinterpret, or supplement it. The evaluation design must therefore measure
both stages rather than reporting only final answer similarity.

\subsection{Sparse Retrieval}

Sparse retrieval represents a query and document using lexical terms. BM25 is a
widely used probabilistic ranking function that rewards term matches while
controlling for document length and repeated term saturation
\cite{robertson2009bm25}. It performs well when the question contains exact
product names, menu labels, error codes, or workflow terminology. It is also
deterministic, inexpensive, and easy to inspect.

Its limitation is vocabulary mismatch. A user may describe ``changing the main
colour'' while a source uses ``site colour palette.'' If the important terms do
not overlap, lexical ranking may miss a semantically relevant passage. The
project therefore treats BM25 as one retrieval mode and as one component of
hybrid search rather than as a universal retriever.

\subsection{Dense Retrieval and Embeddings}

Dense retrieval maps queries and documents into vectors whose proximity
represents semantic similarity. Dense Passage Retrieval demonstrated the value
of learned encoders for open-domain QA \cite{karpukhin2020dpr}. Sentence-BERT
made transformer representations practical for semantic similarity and search
\cite{reimers2019sentencebert}. The implementation uses
\texttt{sentence-transformers/all-MiniLM-L6-v2} as its principal local
384-dimensional embedding model, with a deterministic hash embedding retained
for dependency-light tests.

Dense retrieval improves semantic matching but introduces a strict integrity
requirement: document vectors and query vectors must be produced by the same
embedding model and compatible preprocessing. A silent model change can make
similarity scores meaningless. This constraint motivates immutable knowledge
index versions and a read-only inherited query-embedding selection in the Main
screen.

Approximate nearest-neighbour indexes, including HNSW
\cite{malkov2018hnsw}, allow vector search to scale without comparing every
vector. PostgreSQL with pgVector was selected so vector similarity can coexist
with relational filters, document metadata, user data, jobs, and evaluation
records \cite{pgvector2026,postgresql2026}.

\subsection{Hybrid Retrieval and Reranking}

Lexical and dense retrievers make different errors. Hybrid retrieval combines
their scores after normalization, allowing exact terminology and semantic
similarity to contribute to the final rank. In \SystemName, retrieval
diagnostics retain the raw BM25 score, raw dense score, normalized components,
weights, hybrid score, rank, and selection status. This makes a poor rank
explainable rather than reducing retrieval to an opaque list of chunks.

A reranker operates on a smaller candidate set after first-stage retrieval. It
can improve precision by evaluating the query and candidate together, but it
adds latency and model cost. The system exposes reranking as an optional
capability selected through a saved RAG configuration. A lexical reranker is
available locally, while hosted rerankers can be registered through the AI
Models layer.

\section{Adaptive Retrieval}

Static RAG applies one retrieval strategy to every question. Adaptive-RAG
instead trains a smaller classifier to select among no retrieval, single-step
retrieval, and iterative retrieval based on question complexity
\cite{jeong2024adaptiverag}. Its labels are derived from strategy outcomes and
dataset inductive biases, which supports the least-complex successful strategy
rather than equating surface length with complexity.

\begin{table}[H]
  \centering
  \caption{Relationship between Adaptive-RAG and the implemented routes.}
  \label{tab:related-route-comparison}
  \begin{tabularx}{\textwidth}{p{0.18\textwidth}p{0.22\textwidth}X}
    \toprule
    Complexity & \SystemName{} route & Execution behavior\\
    \midrule
    Simple & L1 Direct & Generation without knowledge retrieval.\\
    Moderate & L2 Simple RAG & One BM25, dense, or hybrid retrieval followed by generation.\\
    Complex & L3 Complex RAG & Deterministic or model-assisted decomposition, retrieval per subquery, deduplication, aggregation, and generation.\\
    Advanced & L4 Advanced RAG & Bounded planner iterations and authorized tools, with L3 fallback. This is a project extension.\\
    \bottomrule
  \end{tabularx}
\end{table}

Confidence matters because the cost of under-routing and over-routing differs.
Under-routing an advanced question can omit evidence, while over-routing a
simple question increases latency and cost. The four-class design therefore
supports probabilities, confidence, top-two margin, and configurable
escalation. Explicit route selection bypasses escalation so evaluation can test
a fixed configuration.

\section{Query Complexity Classification}

BERT introduced deep bidirectional transformer pre-training for language
understanding \cite{devlin2019bert}. DistilBERT compresses BERT through
knowledge distillation to reduce model size and inference cost
\cite{sanh2019distilbert,hinton2015distilling}. It is suitable for a
sequence-classification router because the output labels are finite and
probabilities can be calibrated.

T5 frames every task as text-to-text generation \cite{raffel2020t5}. For
complexity routing, the model receives a question and assigns probability to
candidate label strings. Candidate scoring is preferable to unrestricted
generation because it guarantees a label from the supported vocabulary. The
project provides both DistilBERT and T5-small training paths through Hugging
Face Transformers \cite{wolf2020transformers}. A deterministic heuristic
remains as a dependency-light runtime fallback.

The implemented classifier uses four labels. The \texttt{advanced} label covers
cross-source dependencies, long conditional chains, and questions that may
require bounded tool use. This is an Aragbiz extension of the original
Adaptive-RAG study's no-retrieval, single-step, and iterative choices.

\section{Iterative and Agentic RAG}

Iterative retrieval refines or decomposes a question and performs multiple
searches. L3 in this project uses a constrained form: two to four subqueries are
produced, each is retrieved independently, repeated chunks are deduplicated,
and evidence is aggregated before one generation call. This increases recall
without granting an open-ended agent control over the system.

Agentic approaches interleave model reasoning with actions. ReAct demonstrated
the value of combining reasoning and action trajectories
\cite{yao2023react}. Self-RAG adds retrieval and critique decisions
\cite{asai2024selfrag}, while Corrective RAG uses retrieval evaluation and
corrective actions \cite{yan2024crag}. These methods motivate L4 but do not
justify an unrestricted production agent. The implemented planner chooses one
validated action per iteration from an authorized registry. Iterations, tool
calls, duration, budgets, egress, and cancellation are bounded. Reasons stored
in traces are concise action rationales rather than private chain-of-thought.

\section{Enterprise Benchmarking with WixQA}

Many QA benchmarks use Wikipedia and do not represent support procedures or
enterprise workflows. WixQA releases a fixed help-centre knowledge base and
three aligned QA sources \cite{cohen2025wixqa}:

\begin{table}[H]
  \centering
  \caption{WixQA sources used by the project.}
  \label{tab:wixqa-sources}
  \begin{tabularx}{\textwidth}{p{0.22\textwidth}p{0.17\textwidth}X}
    \toprule
    Source & Paper size & Characteristics\\
    \midrule
    ExpertWritten & 200 & Authentic user questions with expert-authored,
    detailed answers, frequently requiring multiple articles or steps.\\
    Simulated & 200 & Expert-validated QA pairs distilled from user--expert
    support conversations.\\
    Synthetic & 6,222 & LLM-generated article-grounded QA pairs intended for
    broad corpus coverage.\\
    \bottomrule
  \end{tabularx}
\end{table}

The local preprocessing snapshot contains \PreparedCorpusDocuments{}
non-empty normalized knowledge documents and 6,221 Synthetic QAC rows. This is
one fewer than the paper's reported 6,222 Synthetic pairs. The report preserves
both facts: the publication count describes the released benchmark, while
\PreparedCorpusDocuments{} describes the deterministic artifact actually
audited and used by this repository. Reproducibility metadata records hashes
and counts so the difference is visible rather than silently corrected.

\section{RAG Evaluation and Explainability}

RAGAS proposes reference-free and reference-aware measurements for RAG
components, including faithfulness and answer relevance
\cite{es2023ragas}. This project implements RAGAS-style proxies alongside
WixQA-oriented metrics. Retrieval measurements include hit rate, MRR,
precision, recall, and NDCG. Generation measurements include token F1, BLEU,
ROUGE-1, ROUGE-2, judge-based factuality, and context recall. Operational
measurements include route distribution, latency, tokens, cost, failures, and
retrieved-context count.

Metrics do not directly prescribe an engineering change. RAGXplain evaluates
the relationship among the question, retrieved context, candidate answer, and
ground truth, then produces an executive summary and prioritized improvement
protocols \cite{cohen2025ragxplain}. In this project, a selected AI Model
deployment acts as the judge. Artifacts are stored per evaluation run and
loaded automatically into an embedded insights viewer. The recommendations are
mapped to concrete RAG Customizer or Knowledge Base fields, such as
\texttt{top\_k}, reranker, chunking strategy, prompt, and citation policy.

\section{Technology Selection}

\begin{table}[H]
  \centering
  \caption{Technology selection and rationale.}
  \label{tab:technology-selection}
  \begin{tabularx}{\textwidth}{p{0.24\textwidth}p{0.21\textwidth}X}
    \toprule
    Concern & Selection & Rationale\\
    \midrule
    API and validation & FastAPI & Async Python, typed validation, OpenAPI,
    SSE-compatible responses, and direct integration with ML code
    \cite{fastapi2026}.\\
    Studio interface & React & Component-based stateful UI for streaming chat,
    resizable panels, configuration workflows, charts, and modals
    \cite{react2026}.\\
    Relational storage & PostgreSQL & Transactions, constraints, JSON metadata,
    row locking, indexing, and mature operational tooling
    \cite{postgresql2026}.\\
    Vector storage & pgVector & Vector similarity within PostgreSQL, including
    dimension-aware storage and HNSW indexing \cite{pgvector2026}.\\
    Model access & Local adapters and LiteLLM & One provider-neutral gateway for
    local, OpenRouter, OpenAI, Gemini, Ollama, and vLLM inference
    \cite{litellm2026}.\\
    ML development & Hugging Face Transformers & Reproducible loading,
    fine-tuning, tokenization, and export for DistilBERT, T5, MiniLM, and
    FLAN-T5 \cite{wolf2020transformers}.\\
    Migrations & Alembic & Versioned PostgreSQL schema evolution and data
    backfills \cite{alembic2026}.\\
    Authentication & JWT and RBAC & Stateless API authentication using a
    standardized token representation \cite{rfc7519}.\\
    \bottomrule
  \end{tabularx}
\end{table}

The selected stack supports the research requirements while remaining usable
on a Windows development workstation and in Docker. PostgreSQL is mandatory
for paid-model execution and full analytics; bounded JSON repositories remain
available for deterministic tests and offline development.
